\documentclass{article}
%\usepackage{graphicx,amsbsy}
%\newcommand{\bfG}{\mathbf{\Gamma}}
%\newcommand{\bfL}{\mathbf{\Lambda}}
\title{Nonequilibrium Molecular Dynamics Integrators Using Maple}
\author{
Debra J. Searles\thanks{Department of Chemistry, St Lucia, University of
Queensland, Queensland, AUSTRALIA.}\\
Department of Chemistry,\\
University of Queensland,\\
searles@chemistry.uq.edu.au
\and
Dennis J. Isbister\thanks{School of Physics,
University of NSW, University College, ADFA, Canberra ACT 2600,
AUSTRALIA.}\\
School of Physics,\\
University of NSW,\\
d-isbister@adfa.edu.au
\and
Denis J. Evans\thanks{Research School of Chemistry,
Australian National University, Canberra ACT, AUSTRALIA.}\\
Research School of  Chemistry, \\
Australian National University\\
evans@rsc.anu.edu.au
}
\date{~~~}
\begin{document}

\maketitle

%\vspace{0.1in}
\begin{center}
\section*{Abstract}
\end{center}
The Conjugate Pairing Rule states that for Hamiltonian systems to
which a Gaussian or Nos\'e-Hoover thermostat is applied, the spectrum
of Lyapunov exponents should exhibit a simple symmetry: that the sum
of conjugate pairs of exponents should be a constant independent
of the pair index. This symmetry is important for it enables one
to calculate the entropy production and thereby the nonequilibrium
transport coefficient, by summing just two exponents, say the maximal
exponents. However, the Conjugate Pairing Rule has never been proved
for the standard computer simulation algorithm for shear flow namely
the Sllod algorithm.

We use the symbolic algebra program, Maple, to examine the validity
of the Conjugate Pairing Rule for the Sllod algorithm.  We do this by
using the multiple precision facility of Maple to carry out computer
simulations for Sllod to very high accuracy. We also use the
symbolic algebra capability of Maple to prove the Conjugate Pairing Rule
for ideal gases subject to thermostatted shear flow and also for
unthermostatted shear flow of fluids at arbitrary density.
\vspace{3.0in}

%{\it Submitted to IMACS Mathematics and Computers in Simulation: Non-Standard
%Applications of Computer Algebra}, 7$^{th}$ April 1997.

\end{document}
